\def\Ponderation{25}
\def\SigleNumero{STT-1920}
\def\NomCours{Méthodes statistiques}
\def\DateEvaluation{28 avril 2026}
\def\HeureEvaluation{10 h 30 à 13 h 20}
\def\Session{Hiver 2026}
\def\NomEvaluation{Examen 3}

% Ne rien mettre entre les accolades si l'enseignant (ou l'enseignante) est aussi le coordonnateur (ou la coordonnatrice)
\def\Coordonnatrice{}
\def\Coordonnateur{}

% Ne rien mettre entre les accolades s'il y a plus d'une section
\def\Enseignant{Julien Miron}
\def\Enseignante{}

% Ne rien mettre entre les accolades s'il s'agit d'un cours à section unique
% Mettre le nom de chacun des enseignants et enseignantes entre les accolades
% Une case à cocher sera générée pour chacune des sections
\def\SectionA{}
\def\SectionB{}
\def\SectionC{}
\def\SectionD{}


\def\Directives{
\begin{enumerate}
\item Matériel autorisé:
\begin{itemize}
\item Calculatrice scientifique {\bf autorisée par la faculté des sciences et de génie};
\item Une seule feuille aide-m\'emoire \textbf{manuscrite} de format lettre (8.5''$\times$11'') recto-verso.
\end{itemize}
\item Assurez-vous que cette évaluation comporte \numquestions ~exercices répartis sur \numpages{}~pages, pour un total de \numpoints{}~points.\
\item \textbf{Veuillez ne pas dégrafer votre examen.}
\item Vous devez\textbf{ justifier toutes vos réponses} par une démarche claire, sauf pour les questions à choix de réponse.\
\item Pour les questions à choix de réponse et les vrais ou faux, vous pouvez mettre un X, un crochet ($\checkmark$) ou remplir la case.\
\item Vous devez laisser tout ce qui ne servira pas durant l'examen à l'avant de la salle AVANT de rejoindre votre place.\
\item Vous n'avez droit à AUCUN appareil électronique en votre possession durant l'examen.
\item Une fois assis à votre place, vous devez demandez l'autorisation pour vous lever, même si l'examen n'est pas encore commencé.
\item Dans les 5 dernières minutes de l'examen, il sera interdit de se lever avant que TOUTES les copies n'aient été ramassées.
\end{enumerate}
\begin{itemize}
\item[$\rightarrow$] Le non-respect des consignes 7-8-9 entraînera une pénalité de 10\%.
\end{itemize}
}

\newcommand{\affichepoints}{}     % ← AVEC points
% \renewcommand{\affichepoints}{on} % ← SANS points (non vide)

\providecommand{\Gradescope}{}


\documentclass{Evaluation}
\usepackage{multicol,graphicx}

\noprintanswers



\begin{document}

\pageblanche

%%%%%%%%%%%%%%%% QUESTIONS 2 ET 3 %%%%%%%%%%%%%%%%



%%%%%%%%%%%%%%%% QUESTION 1 %%%%%%%%%%%%%%%%
\begin{questions}

\question Un ingénieur souhaite comparer la performance de trois marques de freins à disque pour vélo de route. La variable d'intérêt sera donc la distance de freinage parcourue pour passer de 40\,km/h à 0\,km/h. L'expérience est répétée $n$ fois pour chaque type de frein. Utilisez la table d'analyse de la variance suivante pour répondre aux deux questions.

\begin{verbatim}
                   SS   DDL       MS       F0    valeur-p   
Traitement      45.77     2   22.885   14.329    0.000001
Erreur          91.03    57    1.597
\end{verbatim}

\begin{parts}
\part[3]   Laquelle des affirmations suivantes est vraie?

\begin{checkboxes}
\choice Nous savons que $n=60$.
\correctchoice L'estimateur de la variance est $1{,}597$.
\choice La variable réponse est la distance de freinage et le facteur est \texttt{Erreur}.
\choice Comme il s'agit d'une régression linéaire simple, il n'y a pas de facteur.
\end{checkboxes}

\begin{solution}
B est vraie. $\hat{\sigma}^2 = \text{MS}_{\text{Residuals}} = \text{SS}_{\text{Residuals}}/\text{DDL}_{\text{Residuals}} = 91{,}03/57 \approx 1{,}597 \approx 1{,}60$.

A) $\text{DDL}_{\text{disque}}=I-1=2$, donc $I=3$ marques. $\text{DDL}_{\text{Residuals}}=n-I=57$, donc $n=60$. Avec $I=3$ marques et $n=60$ observations, $k=n/I=20$ répétitions par marque (pas $k=60$).

C) Le facteur est \texttt{disque}, pas \texttt{Residuals}.

D) Il s'agit d'une ANOVA, pas d'une régression.
\end{solution}

\vspace{0.5cm}
\part[3] Quelle conclusion peut-on tirer à partir de cette table d'analyse de la variance, au seuil de $1\%$?

\begin{checkboxes}
\choice La variance de la distance de freinage d'au moins une marque de frein est différente d'au moins une autre marque.
\correctchoice La distance de freinage moyenne d'au moins une marque de frein est différente d'au moins une autre marque.
\choice Toutes les moyennes sont égales.
\choice Toutes les moyennes sont différentes.
\choice Toutes les variances sont différentes.
\end{checkboxes}

\begin{solution}
B est vraie. La valeur-$p=0{,}000001 < 0{,}01 = \alpha$, on rejette $H_0$. Au moins une moyenne de distance de freinage diffère des autres.

A) L'ANOVA compare les moyennes, pas les variances.

C) Il s'agit d'une ANOVA à un seul facteur; il n'y a pas de terme d'interaction.
\end{solution}
 
\end{parts}


%%%%%%%%%%%%%%%% QUESTION 5 %%%%%%%%%%%%%%%%

\vspace{0.5cm}

\question Nous faisons un test pour confronter $H_0: \mu=15$ à $H_1:\mu\neq15$. La règle de décision est de rejeter $H_0$ si $\overline{X}<10$ ou si $\overline{X}>20$. Dans chacun des cas suivants, déterminez si une erreur est commise et déterminez de quel type, le cas échéant. 

\vspace{0.5cm}
\begin{parts}
    \part[2] La vraie valeur de $\mu$ est 16 et nous obtenons $\overline{x}=12$.
    
    \begin{oneparcheckboxes}
        \choice Erreur de type I
        \correctchoice Erreur de type II
        \choice Aucune erreur
    \end{oneparcheckboxes}

\vspace{0.5cm}
        \part[2] La vraie valeur de $\mu$ est 15 et nous obtenons $\overline{x}=12$.
    
    \begin{oneparcheckboxes}
        \choice Erreur de type I
        \choice Erreur de type II
        \correctchoice Aucune erreur
    \end{oneparcheckboxes}

\vspace{0.5cm}
        \part[2] La vraie valeur de $\mu$ est 15 et nous obtenons $\overline{x}=21$.
    
    \begin{oneparcheckboxes}
        \correctchoice Erreur de type I
        \choice Erreur de type II
        \choice Aucune erreur
    \end{oneparcheckboxes}

\vspace{0.5cm}
        \part[2] La vraie valeur de $\mu$ est 18 et nous obtenons $\overline{x}=21$.
    
    \begin{oneparcheckboxes}
        \choice Erreur de type I
        \choice Erreur de type II
        \correctchoice Aucune erreur
    \end{oneparcheckboxes}
\end{parts}


\pageblanche
\question Une équipe en sciences biomédicales veut vérifier s'il existe une différence de réponse inflammatoire moyenne (score quantitatif) selon la concentration d'un agent actif dans une formulation. \footnote{Merci à Claude pour la suggestion de mise en situation.}

Elle compare donc 8 concentrations différentes (10\%, 20\%, 25\%, 30\%, 35\%, 40\%, 45\% et 50\%) en réalisant 4 mesures indépendantes pour chaque concentration.
La table d'analyse de la variance suivante est produite par le logiciel utilisé par l'équipe.
Dans celle-ci, certaines quantités ont été remplacées par \texttt{AAA}, \texttt{BBB} et \texttt{???}. Nous considérons que les postulats du modèle sont respectés.

\begin{verbatim}
ANOVA

                     SS      ddl      MS      F0
traitement       1154.3      ???     ???     AAA
erreur            141.4      ???     ???
Total               BBB      ???
\end{verbatim}

\begin{parts}
\part[3] Que vaut \texttt{AAA}?

\begin{checkboxes}
\choice 164,9
\choice 8,163
\correctchoice 27,989
\choice 1154,3
\choice Il n'est pas possible de trouver la valeur de \texttt{AAA} avec les informations données.
\end{checkboxes}

\begin{solution}
$\texttt{AAA}=F_0=\dfrac{MS_{\text{trt}}}{MS_E}=\dfrac{SS_{\text{trt}}/(a-1)}{SS_E/(N-a)}=\dfrac{1154{,}3/7}{141{,}4/24}=\dfrac{164{,}9}{5{,}892}\approx 27{,}989$.
\end{solution}



\vspace{0.5cm}
\part[3] Que vaut \texttt{BBB}?

\begin{checkboxes}
\choice 8,163
\correctchoice 1295,7
\choice 1012,9
\choice 0,122
\choice Il n'est pas possible de trouver la valeur de \texttt{BBB} avec les informations données.
\end{checkboxes}

\begin{solution}
$\texttt{BBB}=SS_{\text{total}}=SS_{\text{trt}}+SS_E=1154{,}3+141{,}4=1295{,}7$.
\end{solution}
\vspace{0.5cm}
\part[3] Quelle est l'estimation de $\sigma^2$?

\begin{checkboxes}
\choice $141{,}4$
\choice $24$
\correctchoice $5{,}89$
\choice $164{,}9$
\choice Il n'est pas possible de donner une estimation de $\sigma^2$ avec les informations données.
\end{checkboxes}

\begin{solution}
$\hat{\sigma}^2=MS_E=\dfrac{SS_E}{\text{DDL}_E}=\dfrac{141{,}4}{24}=5{,}89$.
\end{solution}
\vspace{0.5cm}
\part[3] Au seuil de 5\%, quelle conclusion pouvons-nous tirer de la table d'analyse de la variance?

\begin{checkboxes}
\choice Il n'est pas possible de tirer de conclusion car nous ne connaissons pas les valeurs remplacées par \texttt{???}.
\choice Il faudrait répéter l'expérience plus de fois pour pouvoir conclure.
\choice Toutes les moyennes sont significativement différentes.
\correctchoice La réponse inflammatoire moyenne d'au moins une concentration est significativement différente de celle d'au moins une autre concentration.
\choice Toutes les moyennes sont significativement égales.
\choice La variance de la réponse inflammatoire d'au moins une concentration est significativement différente de celle d'au moins une autre concentration.
\end{checkboxes}

\begin{solution}
La valeur-$p$ associée au test $F$ est très petite (dans la sortie complète: $4{,}6\times 10^{-10}$), donc $p<0{,}05$. On rejette $H_0$ et on conclut qu'au moins une moyenne de réponse inflammatoire diffère entre les concentrations.
\end{solution}

\end{parts}


\pageblanche
\question[10] Soit $X_1,\hdots, X_{16}$ qui proviennent d'une loi normale d'espérance $\mu$ et de variance $\sigma^2$. \\
On fait le test d'hypothèses pour lequel
$H_0:\mu=5$ et $H_1:\mu>5$ au seuil $\alpha=1\%$.\\
Pour quelles valeurs de $s^2$ rejette-t-on $H_0$ si $\overline{x}=6,301$?

\begin{solution}
    Comme la variance est inconnue et que les données sont normales, on utilise
    la statistique

    \[
    T_0=\frac{\overline{X}-\mu_0}{S/\sqrt{n}}\sim t_{n-1}\quad\text{sous }H_0.
    \]

    Ici, $n=16$, $\mu_0=5$ et $\overline{x}=6{,}301$, donc

    \[
    T_0=\frac{6{,}301-5}{s/4}=\frac{5{,}204}{s}.
    \]

    Au seuil $\alpha=1\%$ pour un test unilatéral à droite,
    on rejette $H_0$ si
    \[
    T_0>t_{0{,}99;15}\approx 2{,}602.
    \]

    Donc
    \[
    \frac{5{,}204}{s}>2{,}602
    \iff s<\frac{5{,}204}{2{,}602}=2
    \iff s^2<4.
    \]

    \textbf{Région de rejet :} on rejette $H_0$ pour $0<s^2<4$.
\end{solution}




\pageblanche
%%%%%%%%%%%%%%%% QUESTION 6 %%%%%%%%%%%%%%%%
\question\label{rls}
Pour des raisons de santé publique, on cherche à savoir s'il est possible d'expliquer le taux maximal d'ozone $O_3$ de la journée par la température à midi donnée en degrés celsius. Afin d'étudier le modèle $Y_i=\beta_0+\beta_1 X_i+ E_i$, en supposant que les postulats de régression linéaire sont respectés, vous disposez des résultats suivants:

\begin{verbatim}
MSE: 160.529 sur 8 degrés de liberté
R-carré:  0.7041,	R-carré ajusté:  0.6671 
F0: 19.03 sur 1 et 8 degrés de liberté,  valeur-p: 0.002403
\end{verbatim}

\begin{align*}
\sum_{i=1}^{10}x_i=219{,}1 ; \quad \sum_{i=1}^{10}x_i^2=5214{,}58;\quad \sum_{i=1}^{10}y_i=1026{,}4; \quad \sum_{i=1}^{10}y_i^2=118446{,}6;\quad \sum_{i=1}^{10}x_iy_i=23650{,}84; \quad s_X^2=46,01.
\end{align*}



\begin{parts}
\part[7] Donnez l'équation de la droite de régression permettant de prédire la concentration d'ozone $O_3$ dans l'air en fonction de la température à midi. Justifiez toutes les étapes de votre démarche.

\begin{solution}
On a $n=10$.

$\bar{x}=\dfrac{219{,}1}{10}=21{,}91$ \quad et \quad $\bar{y}=\dfrac{1026{,}4}{10}=102{,}64$.

$S_{XX} = 5214{,}58 - 10(21{,}91)^2 = 5214{,}58 - 4800{,}48 = 414{,}10$.

$S_{XY} = 23650{,}84 - 10(21{,}91)(102{,}64) = 23650{,}84 - 22488{,}42 = 1162{,}42$.

$\hat{\beta}_1 = \dfrac{S_{XY}}{S_{XX}} = \dfrac{1162{,}42}{414{,}10} \approx 2{,}81$.

$\hat{\beta}_0 = \bar{y} - \hat{\beta}_1\bar{x} = 102{,}64 - 2{,}81 \times 21{,}91 \approx 102{,}64 - 61{,}57 = 41{,}07$.

L'équation de la droite de régression est~: $\hat{y} = 41{,}07 + 2{,}81\,x$.
\end{solution}

\vspace{11cm}

\part[2] Donnez le pourcentage de la variabilité totale de la concentration d'ozone $O_3$ dans l'air qui est expliquée par la température à midi.

\begin{solution}
Le coefficient de détermination $R^2 = 0{,}7041$ (donné dans la sortie). Donc $70{,}41\,\%$ de la variabilité totale de la concentration d'ozone est expliquée par la température à midi.
\end{solution}

\vfill
\pagesuivante{\ref{rls}}


\part[1] Un chercheur aimerait estimer la concentration d'ozone $O_3$ dans l'air lors d'une journée précise. Le midi de cette journée, la température est de $14$ degrés. La prévision ponctuelle de la concentration d'ozone $O_3$ dans l'air en cette journée est environ de $80{,}38$. Choisissez l'intervalle de valeurs ayant $95\%$ de chances de contenir la concentration d'ozone $O_3$ dans l'air en cette journée. 

\begin{checkboxes}
\choice $80{,}38\pm t_{0{,}025,\,n-2}\sqrt{\text{MSE}\bigg[\dfrac{1}{n}+\dfrac{(14-\bar{x})^2}{S_{XX}}\bigg]}$
\correctchoice $80{,}38\pm t_{0{,}025,\,n-2}\sqrt{\text{MSE}\bigg[1+\dfrac{1}{n}+\dfrac{(14-\bar{x})^2}{S_{XX}}\bigg]}$
\end{checkboxes}

\begin{solution}
B est le bon choix. On cherche un \textbf{intervalle de prévision} pour une observation future individuelle (la concentration d'ozone «\,en une journée précise\,»). L'intervalle de prévision comporte le terme $1+\frac{1}{n}+\frac{(x_0-\bar{x})^2}{S_{XX}}$ sous la racine, alors que l'intervalle de confiance pour $E(Y\mid x_0)$ ne comporte que $\frac{1}{n}+\frac{(x_0-\bar{x})^2}{S_{XX}}$.
\end{solution}

\vspace{1cm}

\part[1] Si la température à midi ce jour-là était de 18 degrés, l'intervalle de valeurs calculé à la question précédente serait...

\begin{checkboxes}
\correctchoice plus court
\choice de même longueur
\choice plus long
\end{checkboxes}

\begin{solution}
A est le bon choix. La largeur de l'intervalle dépend de $(x_0-\bar{x})^2/S_{XX}$. On a $\bar{x}\approx 21{,}91$. Avec $x_0=14$~: $(14-21{,}91)^2=62{,}57$. Avec $x_0=18$~: $(18-21{,}91)^2=15{,}29$. Puisque $18$ est plus proche de $\bar{x}$, le terme $(x_0-\bar{x})^2$ est plus petit, donc l'intervalle est \textbf{plus court}.
\end{solution}

\vspace{1cm}

\part[1] Quelle affirmation parmi les suivantes est vraie au seuil $1\%$?

\begin{checkboxes}
\choice La concentration d'ozone $O_3$ dans l'air et la température à midi ne sont pas unis par une relation linéaire significative.
\choice On ne peut pas prédire la concentration d'ozone $O_3$ dans l'air à partir de la température à midi.
\choice La concentration d'ozone $O_3$ dans l'air et la température à midi sont indépendantes.
\correctchoice La relation linéaire entre la concentration d'ozone $O_3$ dans l'air et la température à midi est significative.
\end{checkboxes}

\begin{solution}
D est le bon choix. La valeur-$p = 0{,}002403 < 0{,}01 = \alpha$. On rejette $H_0\!: \beta_1=0$. La relation linéaire entre la concentration d'ozone et la température à midi est significative au seuil de $1\,\%$.
\end{solution}

\end{parts}


\pageblanche
%%%%%%%%%%%%%%%% QUESTION 7 — Test du khi-carré d'ajustement %%%%%%%%%%%%%%%%
\question Une chocolaterie prétend que ses sacs de 120 bonbons contiennent les saveurs suivantes dans les proportions indiquées~: 30\,\% chocolat noir, 25\,\% chocolat au lait, 25\,\% caramel et 20\,\% noisette. Pour vérifier cette affirmation, un contrôleur de qualité prélève un sac au hasard et compte les bonbons de chaque saveur. Les résultats sont présentés ci-dessous.

\begin{center}
\begin{tabular}{l|c|c|c|c|c}
\textbf{Saveur} &\textsc{ } Noir \textsc{ }& \textsc{ } Lait \textsc{ }  & Caramel & Noisette & \textbf{Total}\\
\hline
\textbf{Fréquence observée} ($O_i$) & 42 & 35 & 20 & 23 & 120\\
\hline
&  &  &  &  & \\

&  &  &  &  & \\

\hline
&  &  &  &  & \\

&  &  &  &  & \\
\hline
\end{tabular}
\end{center}

\begin{parts}
\part[2] Formulez les hypothèses nulle et alternative du test du khi-carré d'adéquation.

\begin{solution}
$H_0$~: La distribution des saveurs correspond aux proportions annoncées, c'est-à-dire $p_{\text{noir}}=0{,}30$, $p_{\text{lait}}=0{,}25$, $p_{\text{caramel}}=0{,}25$, $p_{\text{noisette}}=0{,}20$.

$H_1$~: La distribution des saveurs ne correspond pas aux proportions annoncées (au moins une proportion diffère).
\end{solution}

\vspace{2cm}
\part[4] Calculez la statistique de test $D_0$. Vous pouvez utiliser le tableau dans votre démarche.

\begin{solution}
On calcule d'abord les fréquences espérées $E_i = n \times p_i$~:

\begin{center}
\begin{tabular}{l|cccc}
\textbf{Saveur} & Noir & Lait & Caramel & Noisette \\
\hline
$E_i$ & $120\times0{,}30=36$ & $120\times0{,}25=30$ & $120\times0{,}25=30$ & $120\times0{,}20=24$
\end{tabular}
\end{center}

Toutes les fréquences espérées sont $\geq 5$. $\checkmark$

$$D_0 = \sum_{i=1}^{4}\frac{(O_i - E_i)^2}{E_i} = \frac{(42-36)^2}{36}+\frac{(35-30)^2}{30}+\frac{(20-30)^2}{30}+\frac{(23-24)^2}{24}$$
$$= \frac{36}{36}+\frac{25}{30}+\frac{100}{30}+\frac{1}{24} = 1 + 0{,}833 + 3{,}333 + 0{,}042 = 5{,}208$$
\end{solution}
\vspace{8cm}
\part[2] Combien de degrés de liberté possède la loi de référence sous $H_0$?

\begin{solution}
$\nu = k - 1 = 4 - 1 = 3$ degrés de liberté ($k=4$ catégories, aucun paramètre estimé à partir des données).
\end{solution}
\vspace{2cm}
\part[3] Au seuil $\alpha=0{,}05$, quelle est votre conclusion? 

\begin{solution}
$D_0 = 5{,}208 < 7{,}815 = \chi^2_{0{,}05;\,3}$.

On ne rejette pas $H_0$. Au seuil de 5\,\%, il n'y a pas suffisamment de preuves pour affirmer que la distribution des saveurs diffère des proportions annoncées par la chocolaterie.
\end{solution}

\end{parts}

\pageblanche
%%%%%%%%%%%%%%%% QUESTION 9 — Test du khi-carré d'indépendance %%%%%%%%%%%%%%%%
\question Une étude cherche à déterminer s'il existe un lien entre le niveau de stress (faible, modéré, élevé) et la qualité du sommeil (bonne, mauvaise) chez des adultes. Un échantillon de 200 personnes donne les résultats suivants.

\begin{center}
\begin{tabular}{c|cc|c}
 & \textbf{Bonne qualité} & \textbf{Mauvaise qualité} & \textbf{Total}\\
\hline
\textbf{Stress faible} & & &\\
$O_{ij}$ & 50 & 10 & 60\\
$E_{ij}$ & 33 & 27& \\
$D_{ij}$ & 8,758 & 10,704 \\
\hline
\textbf{Stress modéré} & & &\\
$O_{ij}$  & 45 & 35 & 80\\
$E_{ij}$ & 44 & 36\\
$D_{ij}$ & 0,023 & 0,028 \\
\hline
\textbf{Stress Élevé} & & &\\
$O_{ij}$  & 15 & 45 & 60\\
$E_{ij}$ & 33 & 27\\
$D_{ij}$ & 9,818 & 12 \\
\hline
\hline
\textbf{Total oberservés} & 110 & 90 & 200
\end{tabular}
\end{center}

\begin{parts}
\part[2] Formulez les hypothèses nulle et alternative du test d'indépendance.

\begin{solution}
$H_0$~: Le niveau de stress et la qualité du sommeil sont indépendants.

$H_1$~: Le niveau de stress et la qualité du sommeil ne sont pas indépendants (il existe une association).
\end{solution}

\vspace{2cm}
\part[4] Calculez la statistique de test $D_0$. 

\begin{solution}
Les fréquences espérées sous $H_0$ sont $E_{ij}=\dfrac{(\text{total ligne }i)\times(\text{total colonne }j)}{n}$.

\begin{center}
\begin{tabular}{l|cc}
 & Bonne qualité & Mauvaise qualité \\
\hline
Stress faible & $\frac{60\times 110}{200}=33$ & $\frac{60\times 90}{200}=27$ \\[4pt]
Stress modéré & $\frac{80\times 110}{200}=44$ & $\frac{80\times 90}{200}=36$ \\[4pt]
Stress élevé & $\frac{60\times 110}{200}=33$ & $\frac{60\times 90}{200}=27$
\end{tabular}
\end{center}

Toutes les fréquences espérées sont $\geq 5$. $\checkmark$

\begin{align*}
D_0 &= \frac{(50-33)^2}{33}+\frac{(10-27)^2}{27}+\frac{(45-44)^2}{44}+\frac{(35-36)^2}{36}\\
&\quad +\frac{(15-33)^2}{33}+\frac{(45-27)^2}{27}\\[4pt]
&= \frac{289}{33}+\frac{289}{27}+\frac{1}{44}+\frac{1}{36}+\frac{324}{33}+\frac{324}{27}\\[4pt]
&= 8{,}758+10{,}704+0{,}023+0{,}028+9{,}818+12{,}000\\[4pt]
&= 41{,}331
\end{align*}
\end{solution}

\vspace{8cm}
\part[2] Combien de degrés de liberté possède la loi de référence sous $H_0$?

\begin{solution}
$\nu = (r-1)(c-1) = (3-1)(2-1) = 2$ degrés de liberté ($r=3$ niveaux de stress, $c=2$ niveaux de qualité du sommeil).
\end{solution}

\vspace{2cm}
\part[3] Au seuil $\alpha=0{,}01$, quelle est votre conclusion? 

\begin{solution}
$D_0 = 41{,}331 > 9{,}210 = \chi^2_{0{,}01;\,2}$.

On rejette $H_0$. Au seuil de 1\,\%, il y a suffisamment de preuves pour conclure que le niveau de stress et la qualité du sommeil ne sont pas indépendants. Les données suggèrent qu'un stress plus élevé est associé à une mauvaise qualité du sommeil.
\end{solution}

\end{parts}


\pageblanche
%%%%%%%%%%%%%%%% QUESTION ANOVA AJOUTEE %%%%%%%%%%%%%%%%

%%%%%%%%%%%%%%%% TESTS D'HYPOTHÈSES SUR MU — 1 ÉCHANTILLON %%%%%%%%%%%%%%%%

\question

\begin{parts}
\part[2] On teste $H_0:\mu=50$ contre $H_1:\mu\neq 50$ et on obtient une valeur-$p$ de $0{,}03$. Au seuil $\alpha=0{,}05$, quelle est la bonne décision?

\begin{checkboxes}
\choice On ne rejette pas $H_0$.
\correctchoice On rejette $H_0$.
\choice Il est impossible de conclure sans la taille d'échantillon.
\end{checkboxes}

\begin{solution}
Comme $0{,}03<0{,}05$, on rejette $H_0$.
\end{solution}

\vspace{0.5cm}
\part[2] Quelle est l'hypothèse nulle du test $F$ global en analyse de la variance?

\begin{checkboxes}
\correctchoice Toutes les moyennes de groupes sont égales.
\choice Toutes les variances de groupes sont égales.
\choice Toutes les proportions de groupes sont égales.
\end{checkboxes}

\begin{solution}
En ANOVA à un facteur, $H_0$ est $\mu_1=\mu_2=\cdots=\mu_a$.
\end{solution}

\vspace{0.5cm}
\part[2] Dans le modèle $Y=\beta_0+\beta_1X+\varepsilon$, que signifie $\beta_1\neq0$?

\begin{checkboxes}
\correctchoice En moyenne, $Y$ change de $\beta_1$ unités quand $X$ augmente d'une unité.
\choice La variance de $Y$ est positive.
\choice $X$ et $Y$ sont indépendantes.
\choice Quand $X$ augment d'une unité, $Y$ change de $\beta_1$ unités.
\end{checkboxes}

\begin{solution}
$\beta_1$ est la pente: si $\beta_1>0$, la relation linéaire moyenne est croissante.
\end{solution}

\vspace{0.5cm}
\part[2] Pour un test d'ajustement avec $k=5$ catégories, combien y a-t-il de degrés de liberté?

\begin{checkboxes}
\choice 3
\correctchoice 4
\choice 5
\end{checkboxes}

\begin{solution}
Pour un test d'ajustement sans paramètre estimé, $\nu=k-1=5-1=4$.
\end{solution}

\vspace{0.5cm}
\part[2] Un test d'indépendance donne une valeur-$p$ de $0{,}002$ au seuil $5\%$. Quelle conclusion est correcte?

\begin{checkboxes}
\choice On ne rejette pas l'indépendance.
\correctchoice On rejette l'indépendance: les variables sont associées.
\choice Les variables sont causales.
\choice Il nous manque le nombre de degrés de liberté pour pouvoir conclure.
\end{checkboxes}

\begin{solution}
Comme $0{,}002<0{,}05$, on rejette l'hypothèse d'indépendance. Le test indique une association, pas une causalité.
\end{solution}

\end{parts}



%%%%%%%%%%%%%%%% QUESTION BONUS %%%%%%%%%%%%%%%%


\end{questions}
\end{document}
